In Questions 2.1, 2.8 and 2.17, provide a reasonable description of the sample space for each of the random experiments.
There can be more than one acceptable interpretation of each experiment.
Describe any assumptions you make.

\enquestion{2.1}{
    Each of three machined parts is classified as either above or below the target specification for the part.
}

\enanswer{
    The sample space can be described as 
    \[
        S = \set{aaa, aab, aba, baa, abb, bab, bba, bbb},
    \]
    where $a = \text{above}$ and $b = \text{below}$.
}

\enquestion{2.8}{
    The following two questions appear on an employee survey questionnaire. 
    Each answer is chosen from the five-point scale 1 (never), 2, 3, 4, 5 (always).  
    \begin{itemize}
        \item Is the corporation willing to listen to and fairly evaluate new ideas?
        \item How often are my coworkers important in my overall job performance?
    \end{itemize}
}

\enanswer{
    The sample space can be described as 
    \[
        S = \{(1,1), (1,2), \dots, (1,5), (2,1), \dots, (5,5)\}
    \]
    The number of possible outcomes in $S$ equals $n_S = 5 \cdot 5 = 25$.
}

\enquestion{2.17}{
    Calls are repeatedly placed to a busy phone line until a connection is achieved.
}

\enanswer{
    The sample space can be described as 
    \[
        S = \set{c, bc, bbc, bbbc, \dots},
    \]
    where $b = \text{busy}$ and $c = \text{connect}$.
}

\enquestion{2.20}{
    Three events are shown on the Venn diagram in the following figure:

    \begin{center}
        \begin{tikzpicture}
            \draw (0, 0) rectangle (4, 4);

            \draw (2.4, 2.8) circle (0.7cm);
            \draw (1.6, 2.8) circle (0.7cm);
            \draw (2, 1) circle (0.7cm);

            \node (A) at ($(1.6, 2.8)+(135:1cm)$) {\small$A$};
            \node (B) at ($(2.4, 2.8)+(45:1cm)$) {\small$B$};
            \node (A) at ($(2,1)+(315:1cm)$) {\small$C$};

        \end{tikzpicture}
    \end{center}

    Reproduce the figure and shade the region that corresponds to each of the following events.

    \begin{tabular}{*{2}{m{0.45\textwidth}}}
        (a) $A'$ & (b) $(A \cap B) \cup (A \cap B')$ \\
        (c) $(A \cap B) \cup C$ & (d) $A \cap (B \cup C)'$ \\
        (e) $(A \cap B)' \cup C$ &
    \end{tabular}
}

\enanswer{
    \begin{tabular}{*{2}{m{0.45\textwidth}}}
        (a) $A'$ is the region outside circle $A$:
        
        \begin{center}
            \begin{tikzpicture}
                \fill[red!30, even odd rule] (0, 0) rectangle (4, 4) (1.6, 2.8) circle (0.7cm);
                \draw (0, 0) rectangle (4, 4);

                \draw (2.4, 2.8) circle (0.7cm);
                \draw (1.6, 2.8) circle (0.7cm);
                \draw (2, 1) circle (0.7cm);

                \node (A) at ($(1.6, 2.8)+(135:1cm)$) {\small$A$};
                \node (B) at ($(2.4, 2.8)+(45:1cm)$) {\small$B$};
                \node (A) at ($(2,1)+(315:1cm)$) {\small$C$};

            \end{tikzpicture}
        \end{center}
    & (b) $(A \cap B) \cup (A \cap B')$ simplifies to $A$, so the entire circle $A$ is shaded:
    
        \begin{center}
            \begin{tikzpicture}
                \fill[red!30] (1.6, 2.8) circle (0.7cm);
                \draw (0, 0) rectangle (4, 4);

                \draw (2.4, 2.8) circle (0.7cm);
                \draw (1.6, 2.8) circle (0.7cm);
                \draw (2, 1) circle (0.7cm);

                \node (A) at ($(1.6, 2.8)+(135:1cm)$) {\small$A$};
                \node (B) at ($(2.4, 2.8)+(45:1cm)$) {\small$B$};
                \node (A) at ($(2,1)+(315:1cm)$) {\small$C$};

            \end{tikzpicture}
        \end{center}
    \\
    (c) $(A \cap B) \cup C$ is the intersection of $A$ and $B$, plus the entire circle $C$:

        \begin{center}
            \begin{tikzpicture}
                % Fill the intersection (A \cap B) using a clipping path
                \begin{scope}
                    \clip (1.6, 2.8) circle (0.7cm);
                    \fill[red!30] (2.4, 2.8) circle (0.7cm);
                \end{scope}

                % Fill the entire circle C
                \fill[red!30] (2, 1) circle (0.7cm);

                % Draw the bounding box and the circles
                \draw (0, 0) rectangle (4, 4);
                \draw (1.6, 2.8) circle (0.7cm);
                \draw (2.4, 2.8) circle (0.7cm);
                \draw (2, 1) circle (0.7cm);

                % Labels
                \node at ($(1.6, 2.8)+(135:1cm)$) {\small$A$};
                \node at ($(2.4, 2.8)+(45:1cm)$) {\small$B$};
                \node at ($(2,1)+(315:1cm)$) {\small$C$};
            \end{tikzpicture}
        \end{center}
    & (d) $A \cap (B \cup C)'$ is the portion of circle $A$ that is outside $B$ and $C$:

    \begin{center}
        \begin{tikzpicture}
            % Clip to circle A, then shade its complement relative to B and C
            \begin{scope}
                \clip (1.6, 2.8) circle (0.7cm);
                \fill[red!30, even odd rule] (0, 0) rectangle (4, 4) 
                    (2.4, 2.8) circle (0.7cm) 
                    (2, 1) circle (0.7cm);
            \end{scope}

            % Draw the bounding box and the circles
            \draw (0, 0) rectangle (4, 4);
            \draw (1.6, 2.8) circle (0.7cm);
            \draw (2.4, 2.8) circle (0.7cm);
            \draw (2, 1) circle (0.7cm);

            % Labels
            \node at ($(1.6, 2.8)+(135:1cm)$) {\small$A$};
            \node at ($(2.4, 2.8)+(45:1cm)$) {\small$B$};
            \node at ($(2,1)+(315:1cm)$) {\small$C$};
        \end{tikzpicture}
    \end{center}
    \\
    (e) $(A \cap B)' \cup C$ is everything outside the intersection of $A$ and $B$, plus the entire circle $C$:

    \begin{center}
        \begin{tikzpicture}
            % Fill the entire bounding box
            \fill[red!30] (0, 0) rectangle (4, 4);
            
            % "Punch out" (unfill) the intersection of A and B with white
            \begin{scope}
                \clip (1.6, 2.8) circle (0.7cm);
                \fill[white] (2.4, 2.8) circle (0.7cm);
            \end{scope}

            % Draw the bounding box and the circles on top
            \draw (0, 0) rectangle (4, 4);
            \draw (1.6, 2.8) circle (0.7cm);
            \draw (2.4, 2.8) circle (0.7cm);
            \draw (2, 1) circle (0.7cm);

            % Labels
            \node at ($(1.6, 2.8)+(135:1cm)$) {\small$A$};
            \node at ($(2.4, 2.8)+(45:1cm)$) {\small$B$};
            \node at ($(2,1)+(315:1cm)$) {\small$C$};
        \end{tikzpicture}
    \end{center}
    & \\
\end{tabular}

}

\enquestion{2.29}{
    The rise time of a reactor is measured in minutes (and fractions of minutes). 
    Let the sample space be positive, real numbers. 
    Define the events $A$ and $B$ as follows: $A = \set{x \mid x < 65}$ and $B = \set{x \mid x > 45.5}$.
    
    Describe each of the following events:

    \begin{tabular}{*{2}{m{0.45\textwidth}}}
        (a) $A'$ & (b) $B'$ \\
        (c) $A \cap B$ & (d) $A \cup B$
    \end{tabular}
}

\enanswer{
    \begin{tabular}{*{2}{m{0.45\textwidth}}}
        (a) $\set{x \mid x \geq 65}$ & (b) $\set{x \mid 0 < x \leq 45.5}$ \\
        (c) $\set{x \mid 45.5 < x < 65}$ & (d) $\set{x \mid 0 < x < \infty}$
    \end{tabular}
}

\enquestion{2.31}{
    A sample of two printed circuit boards is selected without replacement from a batch.
    Describe the (ordered) sample space for each of the following batches:
    \begin{enumerate}[label=(\alph*)]
        \item The batch contains $90$ boards that are not defective, seven boards with minor defects, and three boards with major defects.
        \item The batch contains $90$ boards that are not defective, eight boards with minor defects, and one board with major defects.
    \end{enumerate}
}

\enanswer{
    \begin{tabular}{@{}l l}
        a. & $S = \set{gg, gm, gb, mg, mm, mb, bg, bm, bb}$, \\
           & where $g = \text{good}$, $m = \text{minor}$ and $b = \text{bad}$ \\
        b. & $S = \set{gg, gm, gb, mg, mm, mb, bg, bm}$
    \end{tabular}
}

\enquestion{2.55}{
    The sample space of a random experiment is $\{a, b, c, d, e\}$ with probabilities $0.2, 0.1, 0.2, 0.3,$ and $0.2$, respectively.
    Let $A$ denote the event $\{a, b, c\}$, and let $B$ denote the event $\{c, d, e\}$.
    Determine the following:

    \begin{tabular}{*{3}{m{.3\textwidth}}}
        (a) $P(A)$ & (b) $P(B)$ & (c) $P(A')$ \\
        (d) $P(A \cup B)$ & (e) $P(A \cap B)$ &
    \end{tabular}
}

\enanswer{
    \begin{enumerate}
        \item $P(A) = P(a) + P(b) + P(c) = 0.2 + 0.1 + 0.2 = 0.5$
        \item $P(B) = P(c) + P(d) + P(e) = 0.2 + 0.3 + 0.2 = 0.7$
        \item $P(A') = 1 - P(A) = 1 - 0.5 = 0.5$
        \item Notice that $A \cup B = \set{a,b,c,d,e} = S$, therefore $P(A \cup B) = 1$.
        \item Notice that $A \cap B = \set{c}$, therefore $P(A \cap B) = P(c) = 0.2$.
    \end{enumerate}
}

\enquestion{2.74}{
    If $P(A) = 0.4$, $P(B) = 0.2$, and $P(A \cap B) = 0.1$, determine the following probabilities:

    \begin{tabular}{*{3}{m{.3\textwidth}}}
        (a) $P(A')$ & (b) $P(A \cup B)$ & (c) $P(A' \cap B)$ \\
        (d) $P(A \cap B')$ & (e) $P[(A \cup B)']$ & (f) $P(A' \cup B)$
    \end{tabular}
}

\enanswer{
    \begin{enumerate}
        \item $P(A') = 1 - P(A) = 1 - 0.4 = 0.6$
        \item $P(A \cup B) = P(A) + P(B) - P(A \cap B) = 0.4 + 0.2 - 0.1 = 0.5$
        \item $P(A' \cap B) = P(B) - P(A \cap B) = 0.2 - 0.1 = 0.1$
        \item $P(A \cap B') = P(A) - P(A \cap B) = 0.4 - 0.1 = 0.3$
        \item $P[(A \cup B)'] = 1 - P(A \cup B) = 1 - 0.5 = 0.5$
        \item $P(A' \cup B) = P(A') + P(B) - P(A' \cap B) = 0.6 + 0.2 - 0.1 = 0.7$. 
    \end{enumerate}
}

\enquestion{2.77}{
    Disks of polycarbonate plastic from a supplier are analyzed for scratch and shock resistance.
    The results from $100$ disks are summarized as follows:

    \begin{center}
        \begin{tabular}{llcc}
            \toprule
                & & \multicolumn{2}{c}{\textbf{Shock Resistance}} \\
            \midrule
                & & High & Low \\
                \textbf{Scratch} & High & $70$ & $10$ \\
                \textbf{Resistance} & Low & $15$ & $5$ \\
            \bottomrule
        \end{tabular}
    \end{center}

    \begin{enumerate}[label=(\alph*)]
        \item If a disk is selected at random, what is the probability that its scratch resistance is high and its shock resistance is high?
        \item If a disk is selected at random, what is the probability that its scratch resistance is high or its shock resistance is high?
        \item Consider the event that a disk has high scratch resistance and the event that a disk has high shock resistance. Are these two events mutually exclusive?
    \end{enumerate}
}

\enanswer{
    \begin{enumerate}[label=(\alph*)]
        \item Of the $100$ disks, there are exactly $70$ with high scratch resistance and high shock resistance (left top cell).
        Therefore the probability of picking such a disk is 
        \[
            \frac{\text{number of favorable outcomes}}{\text{total number of outcomes}} = \frac{70}{100} = 0.7.
        \]
        \item Of the $100$ disks, there are exactly $70 + 10 + 15 = 95$ with high scratch resistance or high shock resistance (all cells in left column or top row).
        Therefore the probability of picking such a disk is 
        \[
            \frac{\text{number of favorable outcomes}}{\text{total number of outcomes}} = \frac{95}{100} = 0.95.
        \]
        \item Let $A$ be the event of high scratch resistance, and $B$ the event of high shock resistance.
        We have that $P(A \cap B) = 0.7 > 0$, hence $A$ and $B$ are not mutually exclusive.
    \end{enumerate}
}

\enquestion{2.87}{
    Disks of polycarbonate plastic from a supplier are analyzed for scratch and shock resistance. 
    The results from 100 disks are summarized as follows:

    \begin{center}
        \begin{tabular}{llcc}
            \toprule
                & & \multicolumn{2}{c}{\textbf{Shock Resistance}} \\
            \midrule
                & & High & Low \\
                \textbf{Scratch} & High & $70$ & $10$ \\
                \textbf{Resistance} & Low & $15$ & $5$ \\
            \bottomrule
        \end{tabular}
    \end{center}

    Let $A$ denote the event that a disk has high shock resistance, and let $B$ denote the event that a disk has high scratch resistance. Determine the following probabilities:

    \begin{tabular}{*{2}{m{0.45\textwidth}}}
        (a) $P(A)$ & (b) $P(B)$ \\
        (c) $P(A \mid B)$ & (d) $P(B \mid A)$
    \end{tabular}
}

\enanswer{
    \begin{enumerate}
        \item Of the $100$ disks, there are exactly $70+15=85$ with high shock resistance (left column).
        Therefore:
        \[
            P(A) = \frac{\text{number of favorable outcomes}}{\text{total number of outcomes}} = \frac{85}{100} = 0.85.
        \]
        \item Of the $100$ disks, there are exactly $70+10=80$ with high shock resistance (top row).
        Therefore: 
        \[
            P(B) = \frac{\text{number of favorable outcomes}}{\text{total number of outcomes}} = \frac{80}{100} = 0.80.
        \]
        \item Of the $80$ disks with high scratch resistance, there are exactly $70$ with high shock resistance (top left cell).
        Therefore:
        \[
            P(A \mid B) = \frac{\text{number of favorable outcomes}}{\text{total number of outcomes}} = \frac{70}{80} = 0.875.
        \]
        \item Of the $85$ disks with high shock resistance, there are exactly $70$ with high scratch resistance (top left cell).
        Therefore:
        \[
            P(B \mid A) = \frac{\text{number of favorable outcomes}}{\text{total number of outcomes}} = \frac{70}{85} = 0.8235.
        \]
    \end{enumerate}
}

\enquestion{2.91}{
    The following table summarizes the analysis of samples of galvanized steel for coating weight and surface roughness:

    \begin{center}
        \begin{tabular}{llcc}
            \toprule
                & & \multicolumn{2}{c}{\textbf{Coating Weight}} \\
            \midrule
                & & High & Low \\
                \textbf{Surface} & High & $36$ & $16$ \\
                \textbf{Roughness} & Low & $114$ & $34$ \\
            \bottomrule
        \end{tabular}
    \end{center}

    \begin{enumerate}[label=(\alph*)]
        \item If the coating weight of a sample is high, what is the probability that the surface roughness is high?
        \item If the surface roughness of a sample is high, what is the probability that the coating weight is high?
        \item If the surface roughness of a sample is low, what is the probability that the coating weight is low?
    \end{enumerate}
}

\enanswer{
    Let us define the events 
    \begin{align*}
        A = \set{\text{Coating weight is high}}, B = \set{\text{Surface roughness is high}}.
    \end{align*}
    The requested probabilities are as follows:
    \begin{enumerate}[label=(\alph*)]
        \item $P(B \mid A) = \frac{P(B \cap A)}{P(A)} = \frac{36}{36+114} = 0.24$.
        \item $P(A \mid B) = \frac{P(A \cap B)}{P(B)} = \frac{36}{36+16} = 0.6923$.
        \item $P(A' \mid B') = \frac{P(A' \cap B)}{P(B')} = \frac{34}{34+114} = 0.2297$.
    \end{enumerate}
}

\enquestion{2.106}{
    Suppose that $P(A \mid B) = 0.25$, $P(A \mid B') = 0.35$, and $P(B) = 0.8$. 
    What is $P(A)$?
}

\enanswer{
    By the complement rule, we have that $P(B') = 1 - P(B) = 1 - 0.8 = 0.2$.
    Using the law of total probability, we get:
    \[
        P(A) = P(A \mid B) \cdot P(B) + P(A \mid B') \cdot P(B') = 0.25 \cdot 0.8 + 0.35 \cdot 0.2 = 0.27.
    \]
}

\enquestion{2.113}{
    A batch of $30$ injection-molded parts contains 6 parts that have suffered excessive shrinkage.
    \begin{enumerate}[label=(\alph*)]
        \item If two parts are selected at random, and without replacement, what is the probability that the second part selected is one with excessive shrinkage?
        \item If three parts are selected at random, and without replacement, what is the probability that the third part selected is one with excessive shrinkage?
    \end{enumerate}
}

\enanswer{
    Let us define for $i=1,2,3$ the events $S_i = \set{\text{the $i$-th part has excessive shrinkage}}$.
    We know that $P(S_1) = \frac{6}{30}=\frac{1}{5}$ and can find the probabilities $P(S_2)$ and $P(S_3)$ as follows:
    \begin{enumerate}[label=(\alph*)]
        \item To compute $P(S_2)$, we can split it in two cases (based on whether or not $S_1$ occurred):
        \begin{itemize}
            \item If $S_1$ occurred, then for the second pick we have $30-1=29$ options of which $6-1=5$ have excessive shrinkage.
            With other words: $P(S_2 | S_1) = \frac{5}{29}$.
            \item If $S_1$ did not occur, then for the second pick we have $30-1=29$ options of which $6-0=6$ have excessive shrinkage.
            With other words: $P(S_2 | S_1') = \frac{6}{29}$.
        \end{itemize}
        Using the law of total probability, we get:
        \[
            P(S_2) = P(S_2 \mid S_1) \cdot P(S_1) + P(S_2 \mid S_1') \cdot P(S_1') = \frac{5}{29} \cdot \frac{1}{5} + \frac{6}{29} \cdot \frac{4}{5} = \frac{29}{145} = 0.2.
        \]
        \item To compute $P(S_3)$, we can split it in four cases (based on whether or not $S_1$ and/or $S_2$ occurred):
        \begin{itemize}
            \item If $S_1$ and $S_2$ occurred, then for the third pick we have $30-2=28$ options of which $6-2=4$ have excessive shrinkage.
            With other words: $P(S_3 | S_1 \cap S_2) = \frac{4}{28}$.
            \item If $S_1$ did not occur, but $S_2$ did, then for the third pick we have $30-2=28$ options of which $6-1=5$ have excessive shrinkage.
            With other words: $P(S_3 | S_1' \cap S_2) = \frac{5}{28}$.
            \item If $S_1$ occurred, but $S_2$ did not, then for the third pick we have $30-2=28$ options of which $6-1=5$ have excessive shrinkage.
            With other words: $P(S_3 | S_1 \cap S_2') = \frac{5}{28}$.
            \item If neither $S_1$ nor $S_2$ occurred, then for the third pick we have $30-2=28$ options of which $6-0=6$ have excessive shrinkage.
            With other words: $P(S_3 | S_1' \cap S_2') = \frac{6}{28}$.
        \end{itemize}
        Using the law of total probability, we get:
        \begin{align*}
            P(S_3)  &= P(S_3 \mid S_1 \cap S_2) \cdot P(S_1 \cap S_2) + P(S_3 \mid S_1' \cap S_2) \cdot P(S_1' \cap S_2) \\
            &\qquad+ P(S_3 \mid S_1 \cap S_2') \cdot P(S_1 \cap S_2') + P(S_3 \mid S_1' \cap S_2') \cdot P(S_1' \cap S_2') \\
                    &= \left(\frac{4}{28} \cdot \frac{6}{30} \cdot \frac{5}{29}\right) + \left(\frac{5}{28} \cdot \frac{24}{30} \cdot \frac{6}{29}\right) + \left(\frac{5}{28} \cdot \frac{6}{30} \cdot \frac{24}{29}\right) + \left(\frac{6}{28} \cdot \frac{24}{30} \cdot \frac{23}{29}\right) \\
                    &= 0.2.
        \end{align*}
    \end{enumerate}
}